\documentclass[10pt,a4paper]{amsart}
\usepackage[margin=19mm]{geometry}
\usepackage{amsmath,amssymb,mathtools}
\usepackage[scr=rsfs,cal=euler]{mathalfa}
\usepackage{xfrac}
\usepackage[unicode,hidelinks,bookmarks=false]{hyperref}
\newcommand{\ie}{i.e.\ }
\newcommand{\cf}{cf.\ }
\newcommand{\resp}{respectively\ }
\newcommand{\Subsection}{\subsection}
\providecommand{\nobreakdash}{\nobreak\hbox{-}}
\setlength{\emergencystretch}{2em}
\multlinegap0pt
\usepackage[shortlabels]{enumitem}
\usepackage{amssymb}
\usepackage{bbm}
\usepackage{xcolor}
\usepackage{tikz}
\usepackage{booktabs}
\usepackage{float}
\usepackage{esint}
\usepackage{multirow}
\newtheorem{lemma}{Lemma}
\newtheorem{cor}{Corollary}
\newcommand{\R}{\mathbb{R}}
\title{Propagation Estimates for the Boson Star Equation}
\author{S\'ebastien Breteaux, J\'er\'emy Faupin, Viviana Grasselli}
\date{Source: arXiv:2512.20718v1 (CC BY 4.0)}
\begin{document}
\maketitle

\noindent\emph{Source and licence.} Propagation Estimates for the Boson Star Equation. Version arXiv:2512.20718v1 (CC BY 4.0), distributed by arXiv under the Creative Commons Attribution 4.0 International licence, http://creativecommons.org/licenses/by/4.0/, as recorded in the arXiv licence metadata for this exact version and shown on the abstract page https://arxiv.org/abs/2512.20718v1. Full licence notice: \texttt{licenses/arxiv-2512.20718-CC-BY-4.0.txt}.\par\smallskip

\noindent\textit{Editorial scope.} Counted unit: the article's cor cor:weighted-sobolev (source0.tex lines 2206--2211). The article's complete original proof of it is delivered with it (source0.tex lines 2213--2310, one \texttt{proof} environment of 98 lines). No mathematics is written, repaired or re-derived: the statement and the proof are the article's own lines, in the article's own notation, and no retained line is substituted or re-flowed.  Retained uncounted as the source-local context the proof needs: lines 2126--2134; lines 2148--2162; lines 2196--2202.  Nothing was omitted from any retained span, so the package carries no omission record.  No retained line contains a citation, so the wrapper carries no bibliography and no bibliography entry.  No retained line contains a figure environment, an image-inclusion command or drawing code, so no image is delivered and the package declares no asset dependency.  Editorial formatting only, in addition: an AMS \texttt{amsart} preamble that restates the article's own declarations and macros used by the retained lines, namely the environments \texttt{lemma}, \texttt{cor} on separate counters, together with the standard packages those lines need.  The retained environments print in this document as Lemma 1, Corollary 1 in order of appearance, whereas the article prints the same items with the numbering of its own full document.  The complete native source of the article as distributed in the arXiv e-print for this exact version is bundled unchanged as \texttt{sources/191591/source0.tex}.
\begin{lemma}
		\label{lemma: expr sqrt}
		Let $0\leq \alpha <2 $. Then there exists $c_\alpha>0$ such that 
		\begin{align*}
		    \langle \nabla \rangle^\alpha =&\, \frac{1}{c_\alpha} \int_0^\infty \big(\, u^{\frac\alpha2-1}- u^{\frac\alpha2} (-\Delta + 1 + u )^{-1}\, \big) \mathrm{d}u ,\\
		     \langle x \rangle^\alpha= &\,\frac{1}{c_\alpha} \int_0^\infty \big(\,u^{\frac\alpha2-1}- u^{\frac\alpha2} (x^2 +1+ u )^{-1}\, \big) \mathrm{d}u,
		\end{align*}
as operators on $\mathcal{S}'(\mathbb{R}^d)$.
	\end{lemma}

	\begin{lemma}
	\label{lemma: resolvent Lp}
	    Let $1\leq p\leq \infty$ and $\lambda,s >0$, then it holds 
	    \begin{equation}\label{eq:bound-resolv_Lp}
	        \|(-\Delta + \lambda )^{-1} \|_{\mathcal{B}(L^p)} \lesssim \lambda^{-1}, \quad  \|\partial_k^m(-\Delta + \lambda )^{-1} \|_{\mathcal{B}(L^p)} \lesssim \lambda^{-m/2}, \quad m =1,2 
	    \end{equation}
	    and 
	    \begin{equation}\label{eq:bound-bessel_Lp}
	        \|\langle \nabla \rangle^{-s}\|_{\mathcal{B}(L^p)} \lesssim 1 , \quad \|\partial_k \langle \nabla \rangle^{-s-1}\|_{\mathcal{B}(L^p)} \lesssim 1
	    \end{equation}
	    for all $k =1,\dots,d$. If $1<p<\infty$, we also have
	    \begin{equation}\label{eq:bound-bessel2_Lp}
	        \|\partial_k \langle \nabla \rangle^{-1}\|_{\mathcal{B}(L^p)} \lesssim 1.
	    \end{equation}
	\end{lemma}

\begin{lemma}
\label{lemma: [x^gamma, grad^s]}
    Let $0\leq \gamma\leq 2$,  $s\geq 0$ and $1\leq p \leq \infty$. Then it holds 
    \begin{equation*}
        \|[\langle x \rangle^\gamma , \langle \nabla \rangle^s]\langle x \rangle^{-\gamma} \langle \nabla \rangle^{-s} \|_{\mathcal{B}(L^p)}\lesssim 1, \quad \|[\langle x \rangle^\gamma , \langle \nabla \rangle^s]\langle \nabla \rangle^{-s}   \langle x \rangle^{-\gamma} \|_{\mathcal{B}(L^p)}\lesssim 1 .
    \end{equation*}
\end{lemma}

\begin{cor}\label{cor:weighted-sobolev}
 Let $0\leq \gamma\leq 2$,  $s\geq 0$ and $1\leq p \leq \infty$. There exist $c, c'>0$ such that 
 \begin{equation*}
    c'\,  \| \langle \nabla \rangle^s \langle x \rangle^\gamma f \|_{L^p} \leq   \| f \|_{H_\gamma^{s,p}} \leq c\,  \| \langle \nabla \rangle^s \langle x \rangle^\gamma f \|_{L^p} . 
 \end{equation*}
\end{cor}

\begin{proof}[Proof of Lemma \ref{lemma: [x^gamma, grad^s]}]
We will prove the lemma for the operator $[\langle x \rangle^\gamma , \langle \nabla \rangle^s]\langle x \rangle^{-\gamma} \langle \nabla \rangle^{-s} $, the other case can be proved analogously. 
We distinguish different cases: 
\begin{itemize}
    \item $s=2$, $\gamma \in \R^+$.

 We compute explicitly the commutator yielding
    \begin{align*}
        [\langle x \rangle^\gamma, \langle \nabla \rangle^2]\langle x \rangle^{-\gamma} \langle \nabla \rangle^{-2} =& [\langle x \rangle^\gamma, -\Delta]\langle x \rangle^{-\gamma} \langle \nabla \rangle^{-2}\\
        = &  \Delta (\langle x \rangle^\gamma) \langle x \rangle^{-\gamma} \langle \nabla \rangle^{-2} +  \langle x \rangle^{-\gamma} \nabla (\langle x \rangle^{\gamma})\cdot \nabla \langle \nabla \rangle^{-2} \\
        & + \nabla ( \langle x \rangle^{\gamma})\cdot \nabla ( \langle x \rangle^{-\gamma})  \langle \nabla \rangle^{-2}, 
    \end{align*}
    which is a sum of bounded operators on $L^p$ thanks to Lemma \ref{lemma: resolvent Lp}.
    
    \item $s \in [0,2)$,  $\gamma = 2$. 
    
    Applying Lemma \ref{lemma: expr sqrt} we obtain the expression
    \begin{align}
    \label{eq: comm x^2}
        [\langle x \rangle^2 , \langle \nabla \rangle^s] \langle x \rangle^{-2} \langle \nabla \rangle^{-s} =& \int_0^\infty u^{s/2}R_{-\Delta} (u) (\Delta (x^2) + 2\nabla(x^2)\cdot \nabla)R_{-\Delta} (u)\,  \langle x \rangle^{-2} \langle \nabla \rangle^{-s} \, \mathrm{d}u
    \end{align}
    where for a non negative selfadjoint operator $A$ we define $R_{A} (u) = (A+1 + u)^{-1}$. For the second term in the integral we can argue as follows: 
    \begin{align*}
        \Big\| \int_0^\infty u^{s/2}R_{-\Delta} (u)& \nabla(x^2)\cdot \nabla R_{-\Delta} (u)\,  \langle x \rangle^{-2} \langle \nabla \rangle^{-s} \, \mathrm{d}u \Big\|_{\mathcal{B}(L^p)} \\
         =& 2 \Big\|\sum_{k=1}^d\int_0^\infty u^{s/2}  R_{-\Delta} (u)^2\, x_k \langle x \rangle^{-2}\partial_k \langle \nabla \rangle^{-s} \, \mathrm{d}u\\
          & \hspace{1cm}+  \int_0^\infty u^{s/2}R_{-\Delta} (u)^2\,x_k \partial_k( \langle x \rangle^{-2}) \langle \nabla \rangle^{-s} \, du\Big\|_{\mathcal{B}(L^p)}\\
          \lesssim & \int_0^\infty  \frac{u^{s/2}}{(1+u)^2} \, \mathrm{d}u \lesssim 1
    \end{align*}
    where we applied Lemma \ref{lemma: resolvent Lp}. Since $\Delta (x^2) = 2d$ the corresponding term in \eqref{eq: comm x^2} is treated applying directly Lemma \ref{lemma: resolvent Lp}. This and the previous inequality, together with expression \eqref{eq: comm x^2} prove the statement for $s \in [0, 2)$ and $\gamma =2$. 
    
    \item $s, \gamma \in [0, 2)$. 
    
    We reason analogously, applying Lemma \ref{lemma: expr sqrt} to obtain the expression
    \begin{align}
        [\langle x \rangle^\gamma , \langle \nabla \rangle^s] =& \int_0^\infty t^{\gamma/2}\,u^{s/2}[ (x^2 + 1+t)^{-1}, (-\Delta + 1 + u)^{-1}]\, \mathrm{d}u\, \mathrm{d}t \nonumber\\
         =& \int_0^\infty t^{\gamma/2}\,u^{s/2} R_{x^2}(t)R_{-\Delta} (u) (\Delta (x^2) + 2\nabla(x^2)\cdot \nabla)R_{-\Delta} (u)R_{x^2}(t)\, \mathrm{d}u\, \mathrm{d}t \label{comp: int commut}. 
    \end{align}
     If $0\leq s<1$ we rewrite the contribution of $\nabla(x^2)\cdot \nabla = 2x\cdot \nabla $ as 
     \begin{align}
        R_{x^2}(t)R_{-\Delta} (u) \, x\cdot \nabla R_{-\Delta} (u)R_{x^2}(t) & \nonumber\\
        =R_{x^2}(t)R_{-\Delta} (u)^2  \big(\,& \sum_{k=1}^d    \partial_k  R_{x^2}(t) x_k + dR_{x^2}(t)-2R_{-\Delta} (u) \Delta R_{x^2}(t) \big). \label{comp: commut x, grad}
    \end{align}
    Using the inequality
    \begin{equation*}
        \|R_{x^2}(t)x_k \langle x \rangle^{-\gamma}\|_{\mathcal{B}(L^p)} \lesssim \begin{cases}
            (1 + t)^{-1} & 1\leq \gamma <2\\
            (1 + t)^{-1/2} & 0 \leq \gamma <1
        \end{cases}
    \end{equation*}
as well as Lemma \ref{lemma: resolvent Lp} to estimate \eqref{comp: commut x, grad} we obtain
    \begin{align}
        \Big\| \Big (\int_0^\infty t^{\gamma/2}\,u^{s/2} R_{x^2}(t)R_{-\Delta} (u) \, \nabla(x^2)\cdot \nabla&\,R_{-\Delta} (u)R_{x^2} (t)\, \mathrm{d}u\, \mathrm{d}t\Big ) \langle x \rangle^{-\gamma} \langle \nabla \rangle^{-s}\Big\|_{\mathcal{B}(L^p)}\nonumber\\
        \lesssim  & \begin{cases}
            \int_0^\infty \frac{t^{\gamma/2}}{(1+t)^2} \mathrm{d}t  \int_0^\infty \frac{u^{s/2}}{(1+u)^{3/2}}  \mathrm{d}u & 1\leq \gamma <2\nonumber\\
            \int _0^\infty \frac{t^{\gamma/2}}{(1+t)^{3/2}} \mathrm{d}t  \int_0^\infty \frac{u^{s/2}}{(1+u)^{3/2}}  \mathrm{d}u & 0\leq \gamma <1\nonumber\\
        \end{cases}\nonumber\\
        \lesssim &\, 1  \label{comp: bound [x^g, grad^s]}
    \end{align}
    since we assumed $s<1$. If $1\leq s <2$ we commute $\partial_k$ to the right in \eqref{comp: commut x, grad} obtaining 
    \begin{align*}
       R_{x^2}(t)R_{-\Delta} (u) \, x\cdot &\nabla R_{-\Delta} (u)R_{x^2}(t) \langle x \rangle^{-\gamma} \langle \nabla \rangle^{-s}\\
        =  R_{x^2}(t) R_{-\Delta} (u)^2 \big( \,& \sum_{k=1}^d R_{x^2}(t) x_k \langle x \rangle^{-\gamma} \partial_k \langle \nabla \rangle^{-s}  +  \partial_k (R_{x^2}(t) x_k \langle x \rangle^{-\gamma} )  \langle \nabla \rangle^{-s}\\
         & + d R_{x^2}(t) \langle x \rangle^{-\gamma} \langle \nabla \rangle^{-s} -2R_{-\Delta} (u) \Delta R_{x^2}(t)  \langle x \rangle^{-\gamma} \langle \nabla \rangle^{-s}\big). 
  \end{align*}
    Since $ \|\partial_k (R_{x^2}(t) x_k \langle x \rangle^{-\gamma} ) \|_{\mathcal{B}(L^p)}\lesssim (1+t)^{-1}$ and applying again Lemma \ref{lemma: resolvent Lp} as before we obtain 
    \begin{align*}
      \Big\| \Big (\int_0^\infty u^{\gamma/2}\,t^{s/2} R_{x^2}(u)R_{-\Delta} (t) \, &\nabla(x^2)\cdot \nabla\,R_{-\Delta} (t)R_{x^2} (u)\, \mathrm{d}u\, \mathrm{d}t\Big ) \langle x \rangle^{-\gamma} \langle \nabla \rangle^{-s}\Big\|_{\mathcal{B}(L^p)}\\
      \lesssim  & \begin{cases}
            \int_0^\infty \frac{t^{\gamma/2}}{(1+t)^2} \mathrm{d}t  \int_0^\infty \frac{u^{s/2}}{(1+u)^{2}}  \mathrm{d}u & 1\leq \gamma <2\\
            \int _0^\infty \frac{t^{\gamma/2}}{(1+t)^{3/2}} \mathrm{d}t  \int_0^\infty \frac{u^{s/2}}{(1+u)^{2}}  \mathrm{d}u & 0\leq \gamma <1\\
        \end{cases}\\
        \lesssim & 1 .
    \end{align*}
    The term in \eqref{comp: int commut} involving $\Delta (x^2)$ can be bounded in an analogous way.
    \end{itemize}
     We have therefore proved the statement holds for $\gamma, s \in [0, 2]$, we now use this fact to prove the case $s \in (2, 4]$. Indeed, we need to prove that 
     \begin{equation*}
         [\langle x \rangle^\gamma , \langle \nabla \rangle^s]\langle x \rangle^{-\gamma} \langle \nabla \rangle^{-s} = \langle x \rangle^\gamma \langle \nabla \rangle^s \langle x \rangle^{-\gamma }\langle \nabla \rangle^{-s} - \mathrm{Id}
     \end{equation*}
     is bounded for $\gamma \in [0,2]$ and $s \in (2, 4]$. We rewrite the first operator in the sum as 
     \begin{align}
     \label{comp: comm s>2}
         \langle x \rangle^\gamma \langle \nabla \rangle^s \langle x \rangle^{-\gamma }\langle \nabla \rangle^{-s} =& \langle x \rangle^\gamma \langle \nabla \rangle^{s-2} \langle x \rangle^{-\gamma }\langle \nabla \rangle^{-s+2} + \langle x \rangle^\gamma \langle \nabla \rangle^{s-2} [-\Delta , \langle x \rangle^{-\gamma}]\langle \nabla \rangle^{-s}
     \end{align}
     where the first operator is bounded, since $\gamma, s-2 \in [0,2]$. For the second operator, let us first consider $ \Delta (\langle x \rangle^{-\gamma}) = c \langle x \rangle^{-\gamma-2} + c' \langle x \rangle^{-\gamma-4} $. Let $ j =2,4$, then we have
     \begin{align*}
          \langle x \rangle^\gamma \langle \nabla \rangle^{s-2} \langle x \rangle^{-\gamma-j} \langle \nabla \rangle^{-s} = \big(\langle x \rangle^\gamma \langle \nabla \rangle^{s-2} \langle x \rangle^{-\gamma}\langle \nabla \rangle^{-s +2}\big)\,\big( \langle \nabla \rangle^{s -2} \langle x \rangle^{-j}  \langle \nabla \rangle^{-s} \big)
     \end{align*}
     where again $\langle x \rangle^\gamma \langle \nabla \rangle^{s-2} \langle x \rangle^{-\gamma}\langle \nabla \rangle^{-s +2}$ is bounded from the previous step and it holds 
     \begin{align*}
         \langle \nabla \rangle^{s -2} \langle x \rangle^{-2}  \langle \nabla \rangle^{-s}  = &  \langle x \rangle^{-2}(\langle x \rangle^{2}\langle \nabla \rangle^{s -2} \langle x \rangle^{-2}  \langle \nabla \rangle^{-s+2}) \langle \nabla \rangle^{-2} ,
     \end{align*}
     which is again a composition of bounded operators on $L^p$, and similarly for 
     \begin{align*}
         \langle \nabla \rangle^{s -2} \langle x \rangle^{-4}  \langle \nabla \rangle^{-s}  = &[\, \langle x \rangle^{-2}(\langle x \rangle^{2}\langle \nabla \rangle^{s -2} \langle x \rangle^{-2}  \langle \nabla \rangle^{-s+2})\,]^2 \langle \nabla \rangle^{-2}.
     \end{align*}
          Hence $ \langle x \rangle^\gamma \langle \nabla \rangle^{s-2} \Delta (\langle x \rangle^{-\gamma})\langle \nabla \rangle^{-s}$ is a bounded operator on $L^p$ and the same can be proved for $ \langle x \rangle^\gamma \langle \nabla \rangle^{s-2} \nabla (\langle x \rangle^{-\gamma})\cdot \nabla\, \langle \nabla \rangle^{-s}$ with similar computations. These facts, combined with the expression given in \eqref{comp: comm s>2} prove the statement for $s \in (2, 4]$. For larger $s$ we reason analogously. 
    \end{proof}

\end{document}
